\documentclass[10pt]{article}
\usepackage[margin=0.85in]{geometry}
\usepackage[T1]{fontenc}
\usepackage{lmodern}
\usepackage{amsmath,amssymb,amsthm}
\usepackage{booktabs}
\usepackage{microtype}
\usepackage[hidelinks]{hyperref}
\newtheorem{theorem}{Theorem}
\newtheorem{proposition}{Proposition}
\newcommand{\Z}{\mathbb Z}
\newcommand{\coker}{\operatorname{coker}}
\title{A Defect--Cokernel Duality for a Restricted Table-of-Marks Map}
\author{Anonymous}
\date{}
\begin{document}
\maketitle

\begin{abstract}
The preceding exact calculations give a restricted $16\times16$ self-mark
matrix $M$ and a three-dimensional C65 saturation defect.  We identify the
embedded defect subgroup inside the C66 cokernel and compute the quotient on
both column and row sides.  The defect subgroup is
$(\mathbb Z/8)\oplus(\mathbb Z/2)^2$, while the quotient has invariant factors
$(1^4,2^8,4^2,12,144)$.  An explicit transpose-side congruence lattice gives
the dual quotient with the same invariants.  Every statement concerns only
the frozen sixteen subgroup types.
\end{abstract}

\section{Question and scope}
Let $G=W(E_6)$ and let $S_1,\ldots,S_{16}$ be the fixed subgroup types from
the preceding source package.  Write
\[
 M_{ij}=|(G/S_j)^{S_i}|,
 \qquad M\in\operatorname{Mat}_{16}(\mathbb Z).
\]
The matrix has determinant $226492416=2^{23}3^3$ and is the exact C64
restricted self-mark map.  C65 supplies three integer kernel vectors
\[
 z_1=S_{10}-S_9,\quad
 z_2=-S_2-S_3-S_5-S_6+S_{11}+S_{12}+S_{13}+S_{14},\quad
 z_3=S_{16}-S_{15}.
\]
Put $C=\Z^{16}/M\Z^{16}$.  C66 computes the ambient Smith invariants,
while C67 computes coordinate-wise inverse denominators.  C68 asks how the
C65 saturation vectors sit inside $C$ and what the transpose-side annihilator
looks like.

The scope firewall is
\[
\texttt{NO\_BAD\_EULER\_OR\_ROOT\_NUMBER}.
\]
We do not claim a full table of marks, a full Burnside ring, arithmetic
resolvents, local fields, Euler factors, root numbers, automorphy, or a
Hilbert--Polya operator.

\section{Defect subgroup in the cokernel}
Let $u_1,u_2,u_3$ be the C65 saturation vectors.  Direct exact multiplication
gives
\[
 Mz_1=8u_1,\qquad Mz_2=2u_2,\qquad Mz_3=2u_3.
\]
Thus the classes $[u_i]$ are torsion classes in $C$.  The least multipliers
are detected by the exact rational matrix $M^{-1}U$, where
$U=[u_1\ u_2\ u_3]$.

\begin{theorem}[Embedded defect and quotient]
Let $D=\langle [u_1],[u_2],[u_3]\rangle\subseteq C$.  Then
\[
 D\cong \mathbb Z/8\oplus\mathbb Z/2\oplus\mathbb Z/2,
\]
and
\[
 C/D\cong (\mathbb Z/2)^8\oplus(\mathbb Z/4)^2
 \oplus\mathbb Z/12\oplus\mathbb Z/144.
\]
Equivalently, the Smith invariants of the augmented matrix $[M\mid U]$ are
\[
 (1,1,1,1,2,2,2,2,2,2,2,2,4,4,12,144).
\]
\end{theorem}

\begin{proof}
The three displayed relations show that the classes have orders dividing
$8,2,2$.  Exact rational inversion shows that the corresponding columns of
$M^{-1}U$ have least denominators $8,2,2$, and enumeration of the $8\cdot2\cdot2$
coefficient classes shows that no nonzero combination is integral.  Hence
$|D|=32$ and its invariant factors are $(2,2,8)$.  Independently, Euclidean
Smith reduction of $[M\mid U]$ gives the displayed diagonal.  Its product is
$7077888=226492416/32$, as required for $C/D$.
\end{proof}

\section{Transpose-side annihilator}
Define the congruence lattice
\[
 A=\{y\in\Z^{16}:z_1^Ty\equiv0\pmod 8,\ z_2^Ty\equiv0\pmod2,\ z_3^Ty\equiv0\pmod2\}.
\]
The three functionals are well-defined on the transpose quotient because
$z_i^TM^T=d_i u_i^T$ with $(d_1,d_2,d_3)=(8,2,2)$.  In the named coordinate
order, the residue rows $([z_1^Te_j]_8,[z_2^Te_j]_2,[z_3^Te_j]_2)$ are
\[
\begin{array}{c|cccccccccccccccc}
j&1&2&3&4&5&6&7&8&9&10&11&12&13&14&15&16\\\hline
\text{residue}&000&010&010&000&010&010&000&000&700&100&010&010&010&010&001&001
\end{array}
\]
so the coordinate types annihilated individually are $S_1,S_4,S_7,S_8$.

\begin{proposition}[Row--column duality]
The lattice $A$ contains $M^T\Z^{16}$ with index $32$, and
\[
 A/M^T\Z^{16}\cong (C/D)^\vee.
\]
Its Smith invariants are
$(1^4,2^8,4^2,12,144)$, identical to those of $C/D$.
\end{proposition}

\begin{proof}
An explicit integral basis $P$ for $A$ has determinant $32$.  Multiplying
$P^{-1}M^T$ gives an integer matrix with Smith diagonal
$(1^4,2^8,4^2,12,144)$.  The congruence equations define the annihilator of
$D$ under the natural pairing between the column and transpose cokernels;
therefore the row quotient is the Pontryagin dual of $C/D$.
\end{proof}

\section{Reproducibility and limitations}
The evidence binds C64, C65, C66, and C67 source bytes by SHA-256.  The
producer and checker use separate exact Smith reductions; a SymPy integer
Smith calculation reproduces both quotient diagonals.  Clean replay passes,
and seventeen hostile semantic mutations are rejected.  No Smith
transformation basis is treated as canonical.  The theorem remains a finite
integral statement on the frozen 16-type support.

\section{Conclusion}
C68 locates the C65 dyadic defect inside the C66 cokernel rather than merely
reporting either group in isolation.  The explicit row congruence lattice
supplies the matching annihilator and makes the column/row duality checkable
without additional representation-theoretic or arithmetic assumptions.

\end{document}
